\documentclass[french,12pt,a4paper]{report}
\usepackage[top=0.8cm,bottom=0.8cm,left=1.5cm,right=1.5cm]{geometry}
\usepackage[T1]{fontenc}
\usepackage[utf8]{inputenc}
\usepackage{setspace}
\usepackage{graphicx}
\usepackage[french,english]{babel}
\usepackage{microtype} % diminue les hyphénations trop fréquentes

\begin{document}
\thispagestyle{empty}
\vspace*{0pt}
\vfill
\begin{small}
\begin{center}
\textbf{Simulations numériques du mouvement Brownien dans les environnements complexes}
\end{center}    
\textbf{Résumé :} 
Il est des phénomènes en physique que l'on retrouve dans de nombreux systèmes, mais dont la description formelle et exacte n'est jamais vraiment atteinte. Le mouvement brownien, processus par lequel les petites particules explorent l'espace qui leur est disponible sous l'effet de l'agitation thermique des molécules du fluide environnant, compte parmi ces notions. Depuis son observation en 1827, puis sa formalisation théorique et sa confirmation expérimentale dans un cas idéal très simple au début des années 1900, la communauté scientifique n'a cessé d'approfondir la compréhension des interactions entre processus browniens et confinement, qu'il soit rigide ou déformable. S'y ajoutent les écoulements visqueux, l'activité propre de la particule dans le cas d'un organisme vivant, ou encore les interactions hydrodynamiques entre particules. Chacun de ces cas d'étude présente toutefois des limitations : le couplage entre mobilité confinée et déformabilité de l'interface reste mal quantifié, l'effet conjoint de l'encombrement et de la proximité d'un mur sur la diffusion est rarement traité dans son ensemble, et l'extension des théories de dispersion classiques aux solutés actifs restait jusqu'ici largement ouverte.
Mes travaux de thèse mobilisent l'outil numérique pour avancer sur ces trois fronts. Je présente une étude théorique et numérique du mouvement brownien à proximité d'une interface molle et fluctuante, l'analyse de la modification de la diffusion par la présence d'autres colloïdes et d'un mur à proximité immédiate du colloïde d'intérêt, et l'extension de la théorie de Taylor-Aris sur la dispersion d'un soluté en écoulement aux solutés actifs. Ces développements permettent d'appréhender plus finement le comportement de systèmes microbiologiques où le confinement, parfois lui-même fluctuant, modifie drastiquement la mobilité, où l'activité est inhérente, et où l'encombrement perturbe la dynamique. Ayant isolé les effets propres à chacun de ces éléments, leur couplage devient envisageable pour étudier, par comparaison, des systèmes réels. 




\textbf{Mots-clés :} Mouvement Brownien, Interactions Hydrodynamiques, Écoulements, Activité, Élasticité\\
%\noindent\rule{\textwidth}{1pt}
\noindent\makebox[\linewidth]{\rule{\textwidth}{0.4pt}}

\selectlanguage{english}
\begin{center}
\textbf{Numerical Simulations of Brownian motion in Complex Environments}
\end{center}
\textbf{Abstract:} 

There are phenomena in physics that recur across many systems, yet whose exact, formal description is never truly achieved. Brownian motion, the process by which small particles explore the space available to them under the thermal agitation of the surrounding fluid's molecules, is one such notion. Since its observation in 1827, followed by its theoretical formalisation and experimental confirmation in a very simple idealised case in the early 1900s, the scientific community has continually deepened its understanding of the interactions between Brownian processes and confinement, whether rigid or deformable. To this are added viscous flows, the particle's own activity in the case of a living organism, and hydrodynamic interactions between particles. Each of these areas of study, however, presents limitations: the coupling between confined mobility and interface deformability remains poorly quantified, the combined effect of crowding and wall proximity on diffusion is rarely treated as a whole, and the extension of classical dispersion theories to active solutes had, until now, remained largely open.

My doctoral work draws on numerical tools to advance on these three fronts. I present a theoretical and numerical study of Brownian motion near a soft, fluctuating interface, an analysis of how diffusion is modified by the presence of other colloids and a nearby wall, and the extension of Taylor-Aris dispersion theory for a solute in flow to active solutes. These developments allow for a finer understanding of microbiological systems where confinement, sometimes itself fluctuating, drastically alters mobility, where activity is inherent, and where crowding is the rule rather than the exception. Having isolated the effects specific to each of these elements, their coupling becomes feasible to study, by comparison, real experimental systems.


\textbf{Keywords:} Brownian Motion, Hydrodynamic Interactions, Flows, Activity, Elasticity \\
\noindent\makebox[\linewidth]{\rule{\textwidth}{0.4pt}}

\vfill
\selectlanguage{french}
\begin{center}
    \textbf{Laboratoire Ondes et Matière d'Aquitaine, Centre National de la Recherche Scientifique}\\
Université de Bordeaux, France.
\end{center}
\end{small}
\vfill
\end{document}
